% =============================================================================
% Paper/main.tex — Climate Change and Agriculture: Evidence from Farmland Transactions
%
% Target journal: Journal of Environmental Economics and Management (JEEM)
%
% Compile: from Paper/ run
%     TEXINPUTS=../Preambles:$TEXINPUTS xelatex -interaction=nonstopmode main.tex
%     BIBINPUTS=..:$BIBINPUTS bibtex main
%     TEXINPUTS=../Preambles:$TEXINPUTS xelatex -interaction=nonstopmode main.tex
%     TEXINPUTS=../Preambles:$TEXINPUTS xelatex -interaction=nonstopmode main.tex
%
% Or use the /compile-latex skill, which performs the 3-pass cycle automatically.
%
% Switch to elsarticle (JEEM submission format) before submission:
%     \documentclass[review,12pt,authoryear]{elsarticle}
% and remove the article-class preamble below.
% =============================================================================

\documentclass[11pt]{article}

% Page geometry — generous margins for working drafts; tighten at submission.
\usepackage[margin=1in]{geometry}
\usepackage{setspace}\doublespacing
\usepackage{microtype}

% Shared palette + math + tikz/booktabs from project preamble. Loads under
% xelatex without inputenc warnings; defines primary-blue, primary-gold, etc.
\input{header}

% Bibliography
\usepackage[round,authoryear]{natbib}
\usepackage[capitalise]{cleveref}

% Table / figure paths — Stata writes to ../Tables and ../Figures relative
% to Paper/. Use \input{../Tables/...} and \includegraphics{../Figures/...}.

\title{Climate Change and Agriculture: \\ {Evidence from Farmland Transactions}}
\author{%
  {[Author Name]} \\
  Korea Rural Economic Institute\thanks{Email: [email]. Thanks to \ldots. [Acknowledgments].}
}
\date{\today}

\begin{document}

\maketitle

\begin{abstract}
\noindent
[Abstract goes here — JEEM target length 150--250 words. Lead with the
research question, identification strategy, headline magnitudes (in policy
units), and one-sentence implication for adaptation policy.]
\medskip

\noindent\textbf{JEL codes:} Q15, Q54, Q57 \\
\noindent\textbf{Keywords:} climate change; agriculture; farmland values;
hedonic; Ricardian; adaptation
\end{abstract}

\bigskip

% =============================================================================
% Body — sections live in Paper/sections/ and are \input{}'d here as authored.
% Empty for now; add as drafted. Keep \input{} calls in narrative order.
% =============================================================================

% \input{sections/01_introduction}
% \input{sections/02_literature}
% \input{sections/03_data}
% \input{sections/04_empirical_strategy}
% \input{sections/05_results}
% \input{sections/06_robustness}
% \input{sections/07_discussion}
% \input{sections/08_conclusion}

\section*{[Working skeleton]}
This is a placeholder skeleton produced when the project was bootstrapped. The
real sections will be authored into \texttt{Paper/sections/} and pulled in via
\verb|\input{}| calls above.

% =============================================================================
% Bibliography
% =============================================================================

\bibliographystyle{aer}     % adjust to JEEM style before submission
\bibliography{../Bibliography_base}

% =============================================================================
% Appendix
% =============================================================================

% \clearpage
% \appendix
% \input{appendix/A_additional_tables}
% \input{appendix/B_data_details}
% \input{appendix/C_robustness}

\end{document}
